\documentclass[11pt,a4paper]{article}
\usepackage[T1]{fontenc}
\usepackage{lmodern,amsmath,amssymb,booktabs,microtype}
\usepackage[margin=25mm]{geometry}
\usepackage[hidelinks]{hyperref}
\newcommand{\interp}{\par\smallskip\noindent\textit{Interpretation.}\ }
\newcommand{\R}{\mathbb{R}}
\title{\small Session 1 --- Economic block\\[0.6em]
\LARGE Optimal Growth as a Discrete-Time Control Problem}
\author{Gauthier Vermandel\\\small \'Ecole polytechnique}
\date{}
\begin{document}
\maketitle

This note formulates the saving decision as a constrained dynamical system.
The aim is to distinguish what is given, what the planner chooses, and what
the equations determine. It then connects the mathematical problem to
\texttt{OptimalGrowth.py}. Basic multivariable calculus and discrete-time
recursions are the only prerequisites.

\section{The general finite-horizon problem}
A control problem combines a state equation with a criterion for ranking
admissible trajectories. It is not defined by an objective alone: the initial
state, feasible controls, time horizon and terminal treatment are also data.

Let $t=0,\ldots,H-1$ index decision periods of length $\Delta>0$ years.
The state $x_t\in X\subseteq\R^n$ records inherited quantities; the control
$a_t\in U_t(x_t)\subseteq\R^m$ is chosen at date $t$. A known path $z_t$
contains exogenous inputs. An admissible trajectory satisfies
\begin{equation}\label{eq:generalstate}
 x_{t+1}=F_t(x_t,a_t;z_t),\qquad
 x_0=\bar x,\qquad a_t\in U_t(x_t),\quad x_t\in X.
\end{equation}
\interp The transition map propagates today's state and decision into
tomorrow's state. The initial condition is data, not a decision. A control
sequence is feasible only if its entire induced state path is admissible.

For a period reward $r_t$, terminal payoff $\Phi$ and annual discount rate
$\rho\geq0$, the problem is
\begin{equation}\label{eq:generalproblem}
 \sup_{(a_0,\ldots,a_{H-1})\ {\rm admissible}}
 J_H,\qquad
 J_H=\sum_{t=0}^{H-1}\beta^t r_t(x_t,a_t;z_t)
       +\beta^H\Phi(x_H),\qquad
 \beta=(1+\rho)^{-\Delta}.
\end{equation}
\interp The sum values rewards during the $H$ decision periods; $\Phi$
values the state left after the final decision. Setting $\Phi=0$ is a
substantive assumption, not an innocuous omission. Multiplying all rewards
and the terminal payoff by the same positive constant leaves the optimizer
unchanged. With a fixed $\Delta$, this covers a common factor $\Delta$
when rewards are recorded as annual flows.

We use a supremum because existence is not automatic. For a finite horizon,
compact feasible control sets, continuous transitions and rewards, and a
nonempty closed feasible set ensure that a maximum is attained.
Differentiability is additionally needed for the first-order conditions below.

\section{Specialization: households, production and aggregation}
The economic question is how much output to consume now and how much to
invest for later. There is one good, one capital stock and an exogenous
population. A social planner internalizes the intertemporal allocation;
there is no government, tax system or decentralized price-setting problem
in this benchmark.

\paragraph{Households and welfare.}
Let $L_t>0$ be population and $C_t>0$ aggregate consumption. A factor
$\kappa>0$ converts model consumption units into the units used in utility.
Constant-relative-risk-aversion (CRRA) utility is
\begin{equation}\label{eq:utility}
 c_t=\frac{\kappa C_t}{L_t},\qquad
 u(c)=
 \begin{cases}
 (c^{1-\gamma}-1)/(1-\gamma),&\gamma>0,\ \gamma\ne1,\\
 \log c,&\gamma=1,
 \end{cases}
 \qquad r_t=L_tu(c_t).
\end{equation}
\interp Welfare is the sum of identical individual utilities.
Marginal utility is $u'(c)=c^{-\gamma}>0$ and decreases with consumption.
In this deterministic model $\gamma$ governs consumption smoothing;
the intertemporal elasticity of substitution is $1/\gamma$. Its usual
risk-aversion interpretation matters only when uncertainty is introduced.

\paragraph{Firms and technology.}
For productivity $A_t>0$, capital $K_t>0$ and $0<\alpha<1$,
\begin{equation}\label{eq:production}
 Y_t=f_t(K_t):=A_tK_t^\alpha L_t^{1-\alpha},
 \qquad f_t'(K_t)=\alpha\frac{Y_t}{K_t}.
\end{equation}
\interp Production uses inherited capital and the available population.
The marginal product $f_t'$ measures the extra annual output produced by
one additional unit of capital; diminishing returns make this marginal
benefit fall as capital increases.

\paragraph{Aggregation and the state transition.}
The saving share $s_t$ allocates output between consumption and investment:
\begin{equation}\label{eq:resources}
 I_t=s_t f_t(K_t),\qquad C_t=(1-s_t)f_t(K_t),
 \qquad Y_t=C_t+I_t.
\end{equation}
\interp Saving is not a separate good: it is the part of output invested
rather than consumed. Once $K_t$, $s_t$, $A_t$ and $L_t$ are known,
$Y_t$, $I_t$ and $C_t$ are derived quantities, not independent controls.

Write $d=(1-\delta)^\Delta$, where $0\leq\delta<1$ is annual depreciation.
The economic state evolves according to
\begin{equation}\label{eq:capital}
 K_{t+1}=F_t(K_t,s_t):=dK_t+\Delta s_t f_t(K_t),
 \qquad K_0=\bar K>0.
\end{equation}
\interp Next period's capital is surviving capital plus investment over
$\Delta$ years. Output, consumption and investment are annual flows; capital
is a stock. This is the course's discrete accounting convention, not the
exact solution of a continuous-time capital equation with within-period
depreciation of new investment.

\paragraph{Exogenous inputs.}
With initial values $A_0,L_0>0$, the module constructs
\begin{equation}\label{eq:drivers}
 A_{t+1}=\frac{A_t}{1-g_Ae^{-\delta_A\Delta(t+1)}},
 \qquad
 L_{t+1}=L_t\left(\frac{L_\infty}{L_t}\right)^{\ell_g}.
\end{equation}
\interp Productivity growth declines over time if $0<g_A<1$ and
$\delta_A>0$; it vanishes in the baseline $g_A=0$. For $0<\ell_g<1$,
the logarithmic population gap to $L_\infty$ contracts by $1-\ell_g$
each step. This is geometric convergence in logs, not the continuous
logistic equation. The baseline $\ell_g=0$ freezes population.
Both $g_A$ and $\ell_g$ are step-specific parameters; changing $\Delta$
requires attention to their calibration.

Substituting the derived quantities gives the complete reduced problem:
\begin{equation}\label{eq:growthproblem}
\boxed{\begin{aligned}
 &\sup_{s_0,\ldots,s_{H-1}}
 \sum_{t=0}^{H-1}\beta^t R_t(K_t,s_t)+\beta^H\Phi(K_H),\\
 &R_t(K,s)=L_tu\!\left(\frac{\kappa(1-s)f_t(K)}{L_t}\right),\\
 &K_{t+1}=dK_t+\Delta s_t f_t(K_t),\quad K_0=\bar K,\\
 &K_t>0,\quad 0\leq s_t<1,\quad t=0,\ldots,H-1 .
\end{aligned}}
\end{equation}
\interp There is one endogenous state, $K_t$, and one control, $s_t$.
The known pair $(A_t,L_t)$ enters through the time-dependent maps $f_t$
and $R_t$. The strict upper bound keeps consumption positive. For a
compact numerical feasible set one can impose $s_t\leq1-\varepsilon_s$
with fixed $\varepsilon_s>0$.

\section{Optimality and economic interpretation}
Dynamic programming makes the trade-off in the objective explicit:
a saving choice sacrifices current utility but changes the value of
the capital inherited tomorrow. The Euler equation is a local expression
of this trade-off, not a substitute for the constraints or terminal data.

Let $V_t(K)$ be optimal remaining welfare, discounted to date $t$.
The Bellman recursion is
\begin{equation}\label{eq:bellman}
 V_H(K)=\Phi(K),\qquad
 V_t(K)=\sup_{0\leq s<1}
 \left\{R_t(K,s)+\beta V_{t+1}\big(F_t(K,s)\big)\right\}.
\end{equation}
\interp The planner compares immediate utility with the discounted
continuation value. With $\Phi=0$, saving in the last period gives no
future reward, so $s_{H-1}=0$. This finite-horizon end effect is why the
numerical planning window must extend beyond the dates we want to plot.

At a differentiable interior optimum, define the shadow value
$\lambda_{t+1}=V_{t+1}'(K_{t+1})$. Differentiating with respect to saving gives
\begin{equation}\label{eq:foc}
 -\kappa u'(c_t)f_t(K_t)
 +\beta\lambda_{t+1}\Delta f_t(K_t)=0,
 \qquad
 \kappa u'(c_t)=\beta\Delta\lambda_{t+1}.
\end{equation}
\interp One extra unit of annual investment costs $\kappa u'(c_t)$
in current aggregate utility and creates $\Delta$ units of capital with
discounted value $\beta\lambda_{t+1}$. At a lower-bound optimum the
derivative of the Bellman objective with respect to $s$ is nonpositive;
at an imposed upper bound it is nonnegative. Equality applies only
to an interior saving choice.

The envelope condition, combined with \eqref{eq:foc} at two consecutive
interior dates, yields
\begin{equation}\label{eq:euler}
 u'(c_t)=\beta u'(c_{t+1})
 \left[d+\Delta f'_{t+1}(K_{t+1})\right],
 \qquad 0\leq t\leq H-2.
\end{equation}
\interp Giving up consumption today earns surviving capital plus its
marginal product tomorrow. Along an optimal interior path, the utility
cost equals the discounted utility benefit. Population does not appear
as a separate multiplier because welfare is population-weighted and
consumption is per capita using the same contemporaneous $L_t$.

For CRRA utility, an equivalent growth-rate expression is
\begin{equation}\label{eq:consumptiongrowth}
 \frac{c_{t+1}}{c_t}
 =\left[\beta\left(d+\Delta\alpha
                  \frac{Y_{t+1}}{K_{t+1}}\right)\right]^{1/\gamma}.
\end{equation}
\interp Consumption rises when the discounted gross return exceeds one.
A larger $\gamma$ dampens the consumption-growth response to a given
return. This equality need not hold at control bounds or across successive
re-optimizations that use different terminal dates.

\section{Numerical implementation and checks}
The module uses direct optimization (also called single shooting):
given a vector of saving rates, it simulates the state recursion and
evaluates welfare. A bound-constrained solver then varies the saving vector.
This differs from solving the Bellman equation on a capital grid.

\paragraph{Dates and array rows.}
Mathematical $t$ labels the beginning of an interval. In the revised module,
row zero stores only initial states; row $j\geq1$ stores the decision and
flows for the interval ending at that row. Precisely,
\begin{equation}\label{eq:codemap}
 \texttt{K[j]}=K_j,\quad
 \texttt{s[j]}=s_{j-1},\quad
 \texttt{(Y,I,C)[j]}=(Y_{j-1},I_{j-1},C_{j-1}).
\end{equation}
\interp Production in row $j$ uses \texttt{K[j-1]}, \texttt{A[j-1]}
and \texttt{L[j-1]}. Its investment produces \texttt{K[j]}.
The initial control and flows remain \texttt{NaN} and are never inputs
to the transition. State and flow curves therefore refer respectively
to dates and intervals ending at those dates.

\paragraph{Objective and horizon.}
The code sets $\kappa=1000$, uses $\Phi=0$, and minimizes negative welfare
with the limited-memory BFGS bound-constrained algorithm (L-BFGS-B).
Its current bounds allow $s=1$ and its utility evaluation floors
consumption at $10^{-12}$. That numerical regularization is not identical
to the strictly positive-consumption problem \eqref{eq:growthproblem};
check that the floor is inactive in a reported solution.

For positive $\rho$ and a threshold $0<\tau<1$, the implemented window is
\begin{equation}\label{eq:window}
 H=1+\min\{m\in\mathbb{N}:\beta^m\leq\tau\},
 \qquad \tau=\texttt{p.toly}.
\end{equation}
\interp The algorithm optimizes $H$ decisions, retains the first one,
moves forward one period and optimizes again. This is a rolling-horizon
policy, not one optimization over the displayed sample. The threshold
limits a discount factor; by itself it is not a bound on welfare
truncation error. The loop requires $\rho>0$, although a fixed-horizon
mathematical problem can also use $\rho=0$.

\paragraph{Calibration, not empirical estimation.}
The baseline isolates capital accumulation by freezing productivity and
population. These values are supplied inputs, not parameters estimated
from observations in this notebook.
\begin{center}
\begin{tabular}{@{}lll@{}}
\toprule
Parameter & Value & Meaning\\
\midrule
$\Delta,\ \bar K$ & $1,\ 3$ & Years per step; initial capital\\
$\alpha,\ \delta$ & $0.30,\ 0.06$ & Capital elasticity; annual depreciation\\
$A_0,\ g_A,\ \delta_A$ & $1,\ 0,\ 0.005$ & Productivity and its growth parameters\\
$L_0,\ \ell_g$ & $1,\ 0$ & Normalized population; convergence speed\\
$L_\infty$ & $10500$ & Inactive population target in the baseline\\
$\rho,\ \gamma$ & $0.015,\ 2$ & Annual discount rate; utility curvature\\
$\tau$ & $10^{-3}$ & Discount-factor threshold\\
\bottomrule
\end{tabular}
\end{center}
If population growth is activated, $L_0$ and $L_\infty$ must share a
consistent unit: a normalized population of one cannot be interpreted
as millions without an explicit conversion.

\paragraph{Verification and experiments.}
First verify positive finite consumption and capital, admissible saving,
the resource identity \eqref{eq:resources}, and the transition
\eqref{eq:capital} on every simulated interval. Then compare interior Euler
residuals, alternative starting guesses and longer planning windows.
A successful solver return is useful but not a proof of global optimality.
Finally, vary depreciation or the capital share one at a time and interpret
the saving response through the return in \eqref{eq:consumptiongrowth}.
Population and productivity experiments should keep units and initial
conditions explicit.

\end{document}
